\clearpage
\VerdierPhysicalPageAnchor{0248}%
\begingroup\small\setlength{\parskip}{0pt}\setlength{\emergencystretch}{1.5em}
\setlength{\abovedisplayskip}{4pt plus 2pt minus 2pt}%
\setlength{\belowdisplayskip}{4pt plus 2pt minus 2pt}%

\VerdierSemanticAnchor{ch4d-4-6-8}
\noindent\textbf{Proposition 4.6.8.}\quad
\textit{Soient $Y$ un complexe de $\mathcal A$ et
$\varepsilon\colon Y\longrightarrow I^{\bullet\bullet}$ une résolution
injective de Cartan--Eilenberg de $Y$. Lorsque $Y$ est un objet de
$\operatorname{comp}^{+}(\mathcal A)$, ou bien lorsque la résolution
$I^{\bullet\bullet}$ est de longueur finie, le complexe simple
$\int_2 I^{\bullet\bullet}$ est un complexe de type injectif et le
morphisme de complexes :}

\VerdierSemanticAnchor{ch4d-4-6-8-morphism}
\[
\textit{$\int_2\varepsilon\colon
Y\longrightarrow\int_2 I^{\bullet\bullet}$}
\]
\textit{est un quasi-isomorphisme.}

\VerdierSemanticAnchor{ch4d-4-6-8-proof-spectral-objects}
En effet, le morphisme $\varepsilon$ induit un morphisme entre le premier
objet spectral de $Y$, à valeurs dans $D(\mathcal A)$ (4.1.5), et l'objet
spectral, à valeurs dans $D(\mathcal A)$, défini par
$I^{\bullet\bullet}$ filtré par son premier degré ((4.1.7), (4.1.8)). Ces
deux objets spectraux sont transformés par le foncteur cohomologique
canonique
$H\colon D(\mathcal A)\longrightarrow\mathcal A$ en objets spectraux à
valeurs dans la catégorie abélienne $\mathcal A$. Il suffit de montrer que
$\varepsilon$ induit un isomorphisme entre ces objets spectraux. Or l'objet
spectral défini par $Y$ est toujours stationnaire (chap.~II, 4.4.2) et,
sous les hypothèses faites, l'objet spectral défini par
$I^{\bullet\bullet}$ est stationnaire. Il suffit donc de montrer que
$\varepsilon$ induit un isomorphisme sur les termes initiaux. Ces termes
initiaux sont d'une part :

\VerdierSemanticAnchor{ch4d-4-6-8-initial-y}
\[
\mathrm I_1^{p,q}=H^q(Y^p),
\]
et d'autre part :

\VerdierSemanticAnchor{ch4d-4-6-8-initial-i}
\[
{}'\mathrm I_1^{p,q}=H^q(I^{p,\bullet}).
\]
De plus, le morphisme induit par $\varepsilon$ sur les termes initiaux est
déduit par passage à la cohomologie du morphisme de complexes :

\VerdierSemanticAnchor{ch4d-4-6-8-epsilon-p}
\[
\varepsilon^p\colon Y^p\longrightarrow I^{p,\bullet}.
\]

\VerdierSemanticAnchor{ch4d-4-6-8-proof-end}
Or ce morphisme est un quasi-isomorphisme, ce qui démontre la proposition.

\VerdierSemanticAnchor{ch4d-4-6-9}
\medskip\noindent\textbf{4.6.9.}\quad
Soient $F\colon D(\mathcal A)\longrightarrow\mathcal B$ un foncteur
cohomologique, $Y$ un objet de $\operatorname{comp}(\mathcal A)$ et
$Y\longrightarrow I^{\bullet\bullet}$ une résolution injective de
Cartan--Eilenberg. Le foncteur $F$ transforme l'objet spectral de
Cartan--Eilenberg correspondant (4.6.7) en un objet spectral à valeurs dans
$\mathcal B$, d'où une suite spectrale appelée \emph{suite spectrale de
Cartan--Eilenberg}. Désignons par $(\mathrm{C.E.})_r^{p,q}$ les termes de
cette suite spectrale (chap.~II, 4.3.3.1). Les termes initiaux de cette
suite spectrale sont (4.1.9.1) :

\VerdierSemanticAnchor{ch4d-4-6-9-initial-terms}
\[
(\mathrm{C.E.})_1^{p,q}
=F^{p+q}\bigl(I^{\bullet,p}[-p]\bigr)
=F^q(I^{\bullet,p}),
\]
\endgroup
